\documentclass[12pt, a4paper]{article}

\usepackage[T1]{fontenc}
\usepackage[utf8]{inputenc}
\usepackage[charter, cal=cmcal]{mathdesign}
\usepackage[scaled=0.92]{helvet}
\usepackage{geometry}
\geometry{a4paper, top=30mm, bottom=30mm, inner=30mm, outer=30mm}

\usepackage{amsmath, amsthm}
\usepackage{amssymb}
\usepackage{xcolor}
\usepackage{tikz}
\usetikzlibrary{arrows.meta, calc}
\usepackage{fancyhdr}
\usepackage{titlesec}
\usepackage{enumitem}
\usepackage{booktabs}
\usepackage{array}
\usepackage{parskip}
\usepackage[colorlinks=true, linkcolor=NavyBlue, urlcolor=NavyBlue, citecolor=NavyBlue]{hyperref}

% ── Colours (Journey 6 — navy blue / gold) ────────────────────────────────────
\definecolor{NavyBlue}{HTML}{1E3A8A}
\definecolor{ModuleColor}{HTML}{2563EB}
\definecolor{PuzzleColor}{HTML}{D97706}
\definecolor{DiscoveryColor}{HTML}{B45309}
\definecolor{ProblemColor}{HTML}{6B21A8}
\definecolor{ReflectColor}{HTML}{0D9488}
\definecolor{LightBlue}{HTML}{EFF6FF}
\definecolor{LightGold}{HTML}{FEF3C7}
\definecolor{LightGreen}{HTML}{ECFDF5}

% ── Custom Environments (No TColorBox for Math style) ─────────────────────────
\newenvironment{puzzlebox}
  {\par\medskip\fillcolorbox{LightGold}{PuzzleColor}{The Puzzle}}
  {\par\medskip}

\newenvironment{discoverybox}
  {\par\medskip\fillcolorbox{LightGreen}{DiscoveryColor}{Discovery}}
  {\par\medskip}

\newenvironment{namedconcept}[1]
  {\par\medskip\fillcolorbox{LightBlue}{NavyBlue}{#1}}
  {\par\medskip}

\newenvironment{reflectionbox}
  {\par\medskip\fillcolorbox{LightBlue}{ReflectColor}{Reflection}}
  {\par\medskip}

\newcommand{\fillcolorbox}[3]{%
  \noindent
  \filldrawbox{#1}{#2}{\textbf{\sffamily\color{white}#3}}\par\nobreak\vspace{4pt}
}

\newcommand{\filldrawbox}[3]{%
  \colorbox{#1}{\parbox{\dimexpr\linewidth-2\fboxsep\relax}{%
    \vspace{4pt}\noindent\hspace{4pt}#3\vspace{4pt}%
  }}%
}

% ── Section formatting ────────────────────────────────────────────────────────
\titleformat{\section}
  {\normalfont\sffamily\Large\bfseries\color{NavyBlue}}
  {\thesection}{1em}{}[\color{NavyBlue!40}{\titlerule[1.5pt]}]
\titleformat{\subsection}
  {\normalfont\sffamily\large\bfseries\color{ModuleColor}}{\thesubsection}{1em}{}

\pagestyle{fancy}
\fancyhf{}
\fancyhead[L]{\small\sffamily\color{NavyBlue}\textit{Discovery Curriculum}}
\fancyhead[R]{\small\sffamily\color{ModuleColor}\textit{Journey 6 $\cdot$ Module 2}}
\fancyfoot[C]{\small\color{NavyBlue}\thepage}
\renewcommand{\headrulewidth}{0.4pt}

\newcommand{\newterm}[1]{\textit{\textcolor{DiscoveryColor}{#1}}}

\begin{document}

\thispagestyle{empty}
\begin{center}
{\small\color{ModuleColor}\textbf{DISCOVERY CURRICULUM}\quad$\cdot$\quad\textbf{JOURNEY 6: ABSTRACTION}}
\vspace{0.6cm}

{\color{NavyBlue}\rule{\linewidth}{2.5pt}}\vspace{0.5cm}

{\Huge\bfseries\color{NavyBlue}Module 2}\\[0.3cm]
{\LARGE\bfseries\color{ModuleColor}Rings, Fields, and Galois Theory}\\[0.4cm]
{\color{PuzzleColor}\rule{0.4\linewidth}{1.5pt}}\\[0.4cm]
{\large\itshape\color{NavyBlue!80}
Addition and multiplication define our arithmetic.\\[0.15cm]
Their abstract properties build rings and fields.\\[0.15cm]
The symmetry of field extensions proves the limits of algebra.}

\vspace{1.4cm}
\begin{tabular}{rl}
\textbf{Journey}         & Journey 6 --- Abstraction \\[4pt]
\textbf{Module}          & 2 of 4 \\[4pt]
\textbf{Prerequisites}   & Module 6.1 (group theory), Journey 3 (vector spaces) \\[4pt]
\textbf{You will discover} & Rings, fields, field extensions, Galois groups, solvability by radicals \\[4pt]
\textbf{Format}          & 8 problems $+$ 2 deep problems $\approx$ 13 pages \\
\end{tabular}
\vspace{1.2cm}

{\color{NavyBlue}\rule{\linewidth}{2.5pt}}
\end{center}

\newpage

\section{Module 6.2 --- Rings, Fields, and Galois Theory}
\label{sec:mod6.2}

\subsection*{Motivation}

We formalized in Module 6.1 the algebra of a single operation using Group Theory.
Now, we expand our study to structures that possess two operations—addition and multiplication—which behave like our familiar arithmetic.
These structures are **Rings** and **Fields**.

By studying the relationships between fields, we resolve one of the most famous problems in mathematical history: the insolvability of the quintic equation.
For centuries, mathematicians sought a general formula to solve fifth-degree polynomial equations using arithmetic and radicals (like the quadratic formula).
By using Galois Theory, we will discover that this algebraic question is ultimately a geometric one.
We will show that the symmetries of a polynomial's roots form a group, and if this group lacks a specific structural property, the roots simply cannot be expressed using radicals.

\subsection{The Quintic Barrier}

We analyze the limits of expressing polynomial roots using radical formulas.

\begin{puzzlebox}
Consider the polynomial equations of degree $n$:
\[
  a_n x^n + a_{n-1} x^{n-1} + \dots + a_1 x + a_0 \;=\; 0.
\]
* For $n=2$, the quadratic formula gives the roots using square roots.
* For $n=3$ and $n=4$, complex formulas (Cardano's and Ferrari's formulas) yield the roots using square and cube roots.
* For $n=5$, no general algebraic formula exists.
\begin{enumerate}[label=(\alph*), leftmargin=*]
  \item What does it mean for a number to be "expressible by radicals"?
        Write down a number that is expressible by radicals, and one that is algebraic but not obviously expressible (e.g., a root of $x^5 - 6x + 3 = 0$).
  \item If we have a field of numbers $F$ (like the rational numbers $\mathbb{Q}$), and we introduce a root of an equation (like $\sqrt{2}$), we create a larger field $E = F(\sqrt{2})$.
        Why must this new set of numbers form a closed algebraic field?
  \item Galois linked the solvability of a polynomial to the structure of the group of permutations of its roots that leave the base field $F$ fixed.
        Explain how the symmetry of roots can prevent us from separating them using radical extensions.
\end{enumerate}
\end{puzzlebox}

\subsection{Rings and Fields}

We define the algebraic structures that generalize addition and multiplication.

\begin{namedconcept}{Rings and Fields}
A \newterm{ring} $(R, +, \cdot)$ is a set $R$ with two binary operations satisfying:
1. $(R, +)$ is an abelian group (identity element $0$).
2. Multiplication is associative: $a(bc) = (ab)c$.
3. Distributive laws: $a(b+c) = ab+ac$, and $(a+b)c = ac+bc$.

A \newterm{field} $(F, +, \cdot)$ is a ring in which multiplication is commutative, there is a multiplicative identity $1 \neq 0$, and every non-zero element $a \in F$ has a multiplicative inverse $a^{-1}$ ($a \cdot a^{-1} = 1$).
\end{namedconcept}

\begin{enumerate}[label=(\alph*), leftmargin=*]
  \item Explain why the set of integers $\mathbb{Z}$ is a ring but not a field.
  \item Show that the rational numbers $\mathbb{Q}$, the real numbers $\mathbb{R}$, and the complex numbers $\mathbb{C}$ are fields.
  \item Let $\mathbb{Z}/p\mathbb{Z}$ be the ring of integers modulo $p$. Show that if $p$ is a prime number, $\mathbb{Z}/p\mathbb{Z}$ is a field (denoted $\mathbb{F}_p$). Why does it fail to be a field if $p$ is composite?
\end{enumerate}

\subsection{Field Extensions as Vector Spaces}

We study how larger fields can be viewed as vector spaces over their subfields.

\begin{namedconcept}{Field Extensions}
If $F$ and $E$ are fields and $F \subseteq E$, we say $E$ is a \newterm{field extension} of $F$, denoted $E/F$.
* **Vector Space structure:** The extension field $E$ can be viewed as a vector space over the base field $F$, where vector addition is field addition and scalar multiplication is field multiplication by elements of $F$.
* **Degree of Extension:** The dimension of $E$ as a vector space over $F$ is called the \newterm{degree} of the extension, denoted $[E:F]$.
\end{namedconcept}

\begin{enumerate}[label=(\alph*), leftmargin=*]
  \item Consider the field extension $\mathbb{Q}(\sqrt{2}) = \{a + b\sqrt{2} : a, b \in \mathbb{Q}\}$.
        Show that $\{1, \sqrt{2}\}$ is a basis for this vector space over $\mathbb{Q}$. What is the degree $[\mathbb{Q}(\sqrt{2}):\mathbb{Q}]$?
  \item If $L/K$ and $K/F$ are finite field extensions, the degrees satisfy the **tower rule**:
  \[
    [L:F] \;=\; [L:K] \cdot [K:F].
  \]
  Calculate the degree of the extension $\mathbb{Q}(\sqrt{2}, \sqrt{3})$ over $\mathbb{Q}$.
\end{enumerate}

\subsection{The Galois Group and Solvability}

We formalize the group of symmetries of a field extension.

\begin{namedconcept}{The Galois Group}
Let $E/F$ be a field extension. The \newterm{Galois group} $\text{Gal}(E/F)$ is the group of all automorphisms $\sigma: E \to E$ that leave the base field $F$ fixed:
\[
  \sigma(x) \;=\; x \qquad \text{for all } x \in F.
\]
The group operation is function composition.
\end{namedconcept}

\begin{enumerate}[label=(\alph*), leftmargin=*]
  \item Let $E = \mathbb{Q}(\sqrt{2})$ and $F = \mathbb{Q}$. Show that any automorphism $\sigma \in \text{Gal}(E/F)$ must map $\sqrt{2}$ to either $\sqrt{2}$ or $-\sqrt{2}$.
        Deduce that $\text{Gal}(\mathbb{Q}(\sqrt{2})/\mathbb{Q}) \cong \mathbb{Z}_2$.
  \item A group $G$ is **solvable** if it can be broken down into nested abelian quotient groups.
        Galois' Theorem states that a polynomial is solvable by radicals if and only if its Galois group (the Galois group of its splitting field) is solvable.
  \item The symmetric group $S_5$ (permutations of 5 elements) has order 120. Its normal subgroup, the alternating group $A_5$ (order 60), is a simple non-abelian group.
        Explain why this makes $S_5$ non-solvable, proving that a general quintic equation cannot be solved by radicals.
\end{enumerate}

\newpage

\section{Part 3 --- Discovery}

\begin{discoverybox}
\textbf{What have we built?}

\textbf{First.} Rings generalize arithmetic with addition and multiplication. Fields are rings where non-zero elements have multiplicative inverses.

\textbf{Second.} An extension field $E/F$ can be viewed as a vector space over the base field $F$. The degree $[E:F]$ is its dimension.

\textbf{Third.} The Galois group $\text{Gal}(E/F)$ consists of field automorphisms of $E$ that leave the base field $F$ fixed. It represents the symmetry group of the roots of a polynomial.

\textbf{Fourth.} A polynomial is solvable by radicals if and only if its Galois group is solvable. The existence of non-solvable symmetric groups ($S_n$ for $n \ge 5$) proves that general polynomial equations of degree 5 or higher cannot be solved using radicals.
\end{discoverybox}

\newpage

\section{Part 4 --- What This Is Called}

\begin{namedconcept}{Field and Extension Parameters}
* **Ring:** An algebraic structure with associative addition and multiplication, where addition forms an abelian group.
* **Field:** A commutative ring with a multiplicative identity where all non-zero elements are invertible.
* **Splitting Field:** The smallest extension field of $F$ in which a polynomial factorizes completely into linear terms.
* **Galois Group:** The group of automorphisms of an extension field that fix the base field.
* **Solvable Group:** A group containing a subnormal series where all quotient factor groups are abelian.
\end{namedconcept}

\newpage

\section{Part 5 --- Solved Exercises}

\begin{enumerate}[label=\textbf{E\arabic*.}]
  \item \textbf{Degree of the Complex Field Extension.}
  
        Prove that the set of complex numbers $\mathbb{C}$ is a field extension of the real numbers $\mathbb{R}$ of degree $[\mathbb{C}:\mathbb{R}] = 2$.
        
        \textbf{Solution:}
        By definition, the complex numbers are defined as:
        \[
          \mathbb{C} \;=\; \{a + b i : a, b \in \mathbb{R}\}, \qquad \text{where } i^2 = -1.
        \]
        We can view $\mathbb{C}$ as a vector space over the field of real numbers $\mathbb{R}$:
        1. **Vector Addition:** For $z_1 = a_1 + b_1 i$ and $z_2 = a_2 + b_2 i$, we have $z_1 + z_2 = (a_1+a_2) + (b_1+b_2) i \in \mathbb{C}$.
        2. **Scalar Multiplication:** For any real scalar $c \in \mathbb{R}$ and $z = a + b i$, we have $c \cdot z = (c a) + (c b) i \in \mathbb{C}$.
        
        Every element $z \in \mathbb{C}$ can be written uniquely as a linear combination of the set $\{1, i\}$ with real coefficients.
        * **Span:** $\text{span}_{\mathbb{R}}(\{1, i\}) = \mathbb{C}$ by definition.
        * **Linear Independence:** Suppose $c_1 \cdot 1 + c_2 \cdot i = 0$ for $c_1, c_2 \in \mathbb{R}$.
          If $c_2 \neq 0$, then $i = -c_1/c_2$, which implies $i \in \mathbb{R}$, a contradiction.
          Therefore, $c_2$ must be 0, which immediately implies $c_1 = 0$.
          Thus, $\{1, i\}$ is linearly independent over $\mathbb{R}$.
        
        Since $\{1, i\}$ spans $\mathbb{C}$ and is linearly independent over $\mathbb{R}$, it forms a basis for $\mathbb{C}$ over $\mathbb{R}$.
        The number of elements in this basis is 2.
        Therefore, the degree of the extension is:
        \[
          [\mathbb{C}:\mathbb{R}] \;=\; \dim_{\mathbb{R}}(\mathbb{C}) \;=\; 2.
        \]
        
  \item \textbf{Galois Group of $\mathbb{Q}(\sqrt{2}, \sqrt{3})/\mathbb{Q}$.}
  
        Prove that the Galois group of the field extension $\mathbb{Q}(\sqrt{2}, \sqrt{3})/\mathbb{Q}$ is isomorphic to the Klein four-group $V_4 \cong \mathbb{Z}_2 \times \mathbb{Z}_2$.
        
        \textbf{Solution:}
        Let $E = \mathbb{Q}(\sqrt{2}, \sqrt{3})$. Any element $\sigma \in \text{Gal}(E/\mathbb{Q})$ is a field automorphism that leaves the rational numbers $\mathbb{Q}$ fixed.
        The field $E$ is generated over $\mathbb{Q}$ by the elements $\sqrt{2}$ and $\sqrt{3}$. Therefore, any automorphism $\sigma$ is completely determined by its values on these two generators: $\sigma(\sqrt{2})$ and $\sigma(\sqrt{3})$.
        
        Since $\sigma$ must preserve algebraic relations:
        \[
          (\sigma(\sqrt{2}))^2 \;=\; \sigma((\sqrt{2})^2) \;=\; \sigma(2) \;=\; 2.
        \]
        This implies $\sigma(\sqrt{2}) \in \{\sqrt{2}, -\sqrt{2}\}$.
        Similarly, $(\sigma(\sqrt{3}))^2 = \sigma(3) = 3 \implies \sigma(\sqrt{3}) \in \{\sqrt{3}, -\sqrt{3}\}$.
        
        There are exactly four possible combinations, giving four automorphisms in the Galois group:
        1. $\sigma_1$ (Identity): $\sigma_1(\sqrt{2}) = \sqrt{2}$, $\sigma_1(\sqrt{3}) = \sqrt{3}$.
        2. $\sigma_2$: $\sigma_2(\sqrt{2}) = -\sqrt{2}$, $\sigma_2(\sqrt{3}) = \sqrt{3}$.
        3. $\sigma_3$: $\sigma_3(\sqrt{2}) = \sqrt{2}$, $\sigma_3(\sqrt{3}) = -\sqrt{3}$.
        4. $\sigma_4$: $\sigma_4(\sqrt{2}) = -\sqrt{2}$, $\sigma_4(\sqrt{3}) = -\sqrt{3}$.
        
        Let us calculate the compositions. For example, evaluate $\sigma_2^2$:
        \[
          \sigma_2^2(\sqrt{2}) \;=\; \sigma_2(-\sqrt{2}) \;=\; -\sigma_2(\sqrt{2}) \;=\; -(-\sqrt{2}) \;=\; \sqrt{2}.
        \]
        \[
          \sigma_2^2(\sqrt{3}) \;=\; \sigma_2(\sqrt{3}) \;=\; \sqrt{3}.
        \]
        Thus, $\sigma_2^2 = \sigma_1$ (the identity).
        By symmetry, we find that all three non-identity elements satisfy:
        \[
          \sigma_2^2 \;=\; \sigma_3^2 \;=\; \sigma_4^2 \;=\; \sigma_1.
        \]
        Furthermore, composition of any two non-identity elements yields the third:
        \[
          \sigma_2 \circ \sigma_3(\sqrt{2}) \;=\; \sigma_2(\sqrt{2}) \;=\; -\sqrt{2}.
        \]
        \[
          \sigma_2 \circ \sigma_3(\sqrt{3}) \;=\; \sigma_2(-\sqrt{3}) \;=\; -\sqrt{3}.
        \]
        which matches $\sigma_4$.
        This multiplication structure is exactly that of the Klein four-group $V_4 \cong \mathbb{Z}_2 \times \mathbb{Z}_2$.
\end{enumerate}

\newpage

\section{Part 6 --- Problems}

\begin{enumerate}[label=\textbf{P\arabic*.}]
  \item \textbf{Ideals and Quotient Rings.}
  
        Let $R$ be a ring. A subset $I \subseteq R$ is an \newterm{ideal} of $R$ if:
        1. $(I, +)$ is a subgroup of $(R, +)$.
        2. For all $r \in R$ and $x \in I$, both $rx \in I$ and $xr \in I$.
        
        \begin{enumerate}[label=(\alph*)]
          \item Show that the set of even integers $2\mathbb{Z}$ is an ideal in the ring of integers $\mathbb{Z}$.
          \item Show that the set of odd integers is not an ideal in $\mathbb{Z}$.
          \item If $I$ is an ideal of $R$, the quotient set $R/I = \{r + I : r \in R\}$ forms a ring under the operations:
          \[
            (a+I) + (b+I) \;=\; (a+b)+I, \qquad (a+I)\cdot(b+I) \;=\; (ab)+I.
          \]
          Explain why the quotient ring $\mathbb{Z}/n\mathbb{Z}$ forms a field if and only if $n$ is a prime number.
        \end{enumerate}

  \item \textbf{Galois Group of a Cubic Polynomial.}
  
        Consider the polynomial $p(x) = x^3 - 2$ over the field of rational numbers $\mathbb{Q}$.
        \begin{enumerate}[label=(\alph*)]
          \item Write down the three roots of $p(x)$ in the complex plane, using the real cube root $\sqrt[3]{2}$ and the primitive cube root of unity $\omega = e^{2\pi i / 3}$.
          \item The splitting field of $p(x)$ over $\mathbb{Q}$ is the smallest field containing all three roots: $E = \mathbb{Q}(\sqrt[3]{2}, \omega)$.
                Show that the degree of this extension is $[E:\mathbb{Q}] = 6$.
          \item Show that the Galois group $\text{Gal}(E/\mathbb{Q})$ has order 6 and is isomorphic to the symmetric group $S_3$ (the symmetries of the triangle from Module 6.1).
          \item Is the group $S_3$ solvable? Is $x^3 - 2 = 0$ solvable by radicals?
        \end{enumerate}
\end{enumerate}

\newpage

\section{Part 7 --- Deep Problems}

\begin{enumerate}[label=\textbf{D\arabic*.}]
  \item \textbf{The Fundamental Theorem of Galois Theory.}
  
        Let $E/F$ be a finite Galois extension with Galois group $G = \text{Gal}(E/F)$.
        The Fundamental Theorem of Galois Theory establishes a bijective correspondence between:
        * The intermediate fields $K$ satisfying $F \subseteq K \subseteq E$.
        * The subgroups $H$ of the Galois group $G$.
        
        The bijection maps a subgroup $H \le G$ to its **fixed field**:
        \[
          E^H \;=\; \{x \in E : \sigma(x) = x \text{ for all } \sigma \in H\}.
        \]
        and maps an intermediate field $K$ to the subgroup fixing it:
        \[
          \text{Gal}(E/K) \;=\; \{\sigma \in G : \sigma(x) = x \text{ for all } x \in K\}.
        \]
        
        \begin{enumerate}[label=(\alph*)]
          \item For the extension $E = \mathbb{Q}(\sqrt{2}, \sqrt{3})$ over $F = \mathbb{Q}$, the Galois group is $G = \{e, \sigma_2, \sigma_3, \sigma_4\} \cong \mathbb{Z}_2 \times \mathbb{Z}_2$.
                List all subgroups of $G$.
          \item For each subgroup, calculate its corresponding fixed field, identifying all intermediate subfields of $E$.
          \item Show that a subgroup $H$ is a normal subgroup of $G$ if and only if its fixed field $E^H$ is a Galois extension of the base field $F$.
        \end{enumerate}

  \item \textbf{The Unsolvable Quintic.}
  
        We prove that the polynomial $p(x) = x^5 - 6x + 3$ is not solvable by radicals over the rationals $\mathbb{Q}$.
        \begin{enumerate}[label=(\alph*)]
          \item Use **Eisenstein's Criterion** with $p=3$ to prove that $p(x)$ is irreducible over $\mathbb{Q}$.
                Deduce that the splitting field $E$ contains a subfield isomorphic to $\mathbb{Q}(\alpha)$, where $\alpha$ is a root, and that 5 divides the degree $[E:\mathbb{Q}]$.
          \item By Cauchy's Theorem, since 5 divides the order of the Galois group $G = \text{Gal}(E/\mathbb{Q})$, the group $G$ must contain an element of order 5.
                Show that the only elements of order 5 in the symmetric group $S_5$ are 5-cycles.
          \item Use calculus (the derivative $p'(x)$ and the intermediate value theorem) to prove that $p(x)$ has exactly three real roots and two complex conjugate roots.
          \item Explain why the automorphism of complex conjugation ($\sigma(z) = \bar{z}$) acts as a transposition (swapping the two complex roots while leaving the three real roots fixed).
          \item A known theorem in group theory states that any subgroup of $S_5$ that contains both a 5-cycle and a transposition is the entire group $S_5$.
                Deduce that $\text{Gal}(E/\mathbb{Q}) \cong S_5$. Prove that the roots of $x^5 - 6x + 3 = 0$ cannot be expressed using radicals.
        \end{enumerate}
\end{enumerate}

\newpage

\section{Part 8 --- Reflection}

\begin{reflectionbox}
\textbf{Conceptual Reflection}

Consider the following questions:
\begin{enumerate}
  \item **On the Geometric Nature of Roots:** We saw that the unsolvability of the quintic is not due to a lack of algebraic creativity, but is an structural constraint of the symmetric group $S_5$. How does this shift our view of algebra from "solving equations" to "studying symmetries"?
  \item **On Galois' Tragic Fate:** Evariste Galois wrote down these foundational ideas the night before he was killed in a political duel at age 20. Why do some of the deepest breakthroughs in abstract mathematics emerge from young, turbulent minds rather than established academic consensus?
  \item **On the Boundaries of Solvability:** The roots of $x^5 - 6x + 3 = 0$ exist in the complex plane, and we can calculate them to arbitrary decimal precision. Why does algebra consider them "unsolvable" simply because they cannot be written with radicals? Is a radical extension physically superior to other functions?
\end{enumerate}
\end{reflectionbox}

\medskip
\noindent\textbf{Looking ahead.}
Module~6.3 introduces **Categories, Functors, and Natural Transformations**. We will elevate our abstraction to the highest level, generalizing not just the elements of a set, but the relationships between different mathematical structures, building the ultimate language of modern mathematics and functional programming.

\vspace{1.4cm}
\begin{center}
{\color{NavyBlue}\rule{0.5\linewidth}{0.8pt}}\\[6pt]
{\small\color{NavyBlue}\textit{End of Journey 6, Module 6.2 --- Rings, Fields, and Galois Theory}}\\[4pt]
{\small\color{NavyBlue}Continue to: \textbf{Module 6.3 --- Categories, Functors, and Natural Transformations}}\\[6pt]
{\color{NavyBlue}\rule{0.5\linewidth}{0.8pt}}
\end{center}

\end{document}
